\chapter*{المقدّمة العامة}
\addcontentsline{toc}{chapter}{المقدّمة العامة}

تُعدّ الأسرة أوّل مؤسّسة اجتماعية يتفاعل داخلها الفرد، ومن خلالها تبدأ عملية التنشئة الاجتماعية
التي تُسهم في بناء شخصيته وتشكيل منظومته القيمية والسلوكية، وتتحدّد فيها أنماط علاقته بذاته
وبالآخرين. ويؤكّد حامد عبد السلام زهران أنّ الأسرة تمثّل الجماعة الاجتماعية الأولى المسؤولة عن
التنشئة، إذ تؤدّي هذا الدور من خلال ما تمارسه من أساليب التوجيه والضبط والثواب والعقاب. ويزداد
تأثير هذه الممارسات خلال مرحلة المراهقة، التي تتميّز بتحوّلات نفسية واجتماعية ومعرفية متسارعة
تجعل المراهق أكثر وعيًا بخبراته الأسرية وأكثر قدرة على تفسيرها وإضفاء المعنى عليها. ومن هذا
المنطلق اهتمّت الدراسة الحالية ببحث العلاقة بين إدراك المراهق لأسلوب المعاملة الوالدية المتسلط
وبين تقديره لذاته، انطلاقًا من قناعة مفادها أنّ فهم الخبرة الأسرية لا يتحقّق بوصف الممارسات
الوالدية في حدّ ذاتها، وإنّما بالكيفية التي يدرك بها المراهق تلك الممارسات ويمنحها معنى.

\medskip
\noindent\textbf{إشكالية الدراسة وتساؤلاتها.}\ تنطلق الدراسة من ملاحظة مؤدّاها أنّ معظم البحوث
التي تناولت العلاقة بين أساليب المعاملة الوالدية وتقدير الذات اعتمدت المقاربة الكمّية وركّزت على
قياس هذه العلاقة إحصائيًا، في حين ظلّت الخبرة الذاتية للمراهق -- أي الكيفية التي يدرك بها
الممارسات الوالدية ويفسّرها -- أقلّ تناولًا، خصوصًا في البيئة المغربية. ومن هنا برزت \emph{فجوة
بحثية} استدعت دراسة كيفية تنطلق من روايات المراهقين أنفسهم. وقد تبلورت الإشكالية في سؤال مركزي
هو: \emph{كيف يدرك المراهقون (15--17 سنة) أسلوب معاملة والديهم المتسلط، وكيف يفسّرون أثر هذا
الإدراك في بناء تقديرهم لذواتهم؟} وتفرّعت عنه ثلاثة أسئلة تتعلّق على التوالي بمظاهر إدراك التسلط،
وبوصف المراهقين لأبعاد تقدير ذاتهم الخمسة، وبتفسيرهم للعلاقة بين المتغيّرين وأوجه التشابه
والاختلاف فيها بين مراهق وآخر.

\medskip
\noindent\textbf{الإطار النظري المعتمد.}\ استندت الدراسة في تناول المفهوم الأول إلى تصنيف ديانا
بومريند لأنماط المعاملة الوالدية، وإلى تطوير ماكوبي ومارتن له، وإلى تصوّر باربر الذي ميّز بين
الضبط السلوكي والتحكّم النفسي، وأكّد أنّ أثر الممارسات الوالدية يرتبط بالطريقة التي يدركها بها
الأبناء لا بوجودها فحسب. وفي ضوء ذلك حُدّدت ستّة أبعاد لإدراك التسلط هي: الضبط والصرامة، والعقاب،
وغياب الحوار والتفسير، والبرود العاطفي وقلّة التقبّل، والرقابة والتدخّل، وكبت الاستقلالية والرأي.
أمّا مفهوم تقدير الذات فاستند إلى تصوّرات روزنبرغ وكوبر سميث وهارتر وماسلو التي تُبرزه بناءً
نفسيًا متعدّد الأبعاد، فاعتُمدت خمسة أبعاد هي: الشخصي والأسري والاجتماعي والأكاديمي والجسمي. وقد
شكّلت هذه الأبعاد -- مجتمعةً -- الأساس الذي بُنيت عليه شبكة المقابلة وعملية الترميز والتحليل، بما
حقّق الاتّساق بين الإطار النظري وأداة الدراسة وطريقة معالجة المعطيات.

\medskip
\noindent\textbf{الدراسات السابقة.}\ أظهرت مراجعة الأدبيات اتّفاقًا عامًّا على أهمّية العلاقة بين
أساليب المعاملة الوالدية وتقدير الذات، مع تباين في نتائجها. فعلى المستوى العربي بيّنت دراسة أيت
مولود وإكردوشن بعلي (2018) ودراسة شكراوي (2025) وجود علاقة بين الأساليب الوالدية وتقدير الذات، إذ
يرتبط الضبط والعقاب وضعف الدفء بانخفاض تقدير الذات. وعلى المستوى الغربي بيّنت دراسة لامبورن
وزملائه (1991) -- على نحو 4100 مراهق -- أنّ المصنّفين لوالديهم ضمن النمط المتسلط يسجّلون أدنى
مستويات تقدير الذات، بينما جاءت الدراسة النيبالية (2023) بنتيجة مغايرة (ارتباط إيجابي ضعيف)
أرجعها أصحابها إلى خصوصية السياق الثقافي. وقد أبرز هذا التباين الحاجة إلى فهمٍ أعمق لا إلى دراسة
كمّية إضافية.

\medskip
\noindent\textbf{المنهج والعيّنة والأداة.}\ اعتمدت الدراسة المنهج الكيفي الاستكشافي، لأنّ سؤالها
المركزي يبدأ بـ«كيف»، ولأنّ هدفها فهم المعاني لا قياس الارتباطات. واختيرت عيّنة قصدية متجانسة
مكوّنة من أحد عشر مراهقًا ومراهقة (خمسة ذكور وستّ إناث) بإحدى الثانويات التأهيلية بمدينة تطوان،
بعد فرزٍ قبلي شمل ثلاثين تلميذًا وتلميذة للتحقّق من إدراكهم أسلوبًا والديًّا متسلطًا. وجُمعت
البيانات عبر المقابلة شبه الموجّهة المبنية على محورَي الدراسة، ثمّ عولجت باستخدام تحليل المضمون
الموجَّه عبر مراحل التفريغ والقراءة المتكرّرة والترميز والتصنيف وحساب التكرارات واستخلاص الاتّجاه
العام وتحديد الحالات الخاصّة.

\medskip
\noindent\textbf{أبرز النتائج.}\ أظهرت الدراسة أنّ المراهقين لا يدركون التسلط من خلال العقاب أو
الضبط وحدهما، بل من خلال منظومة من الممارسات المتداخلة (غياب الحوار، الانفراد بالقرار، الرقابة،
الشعور بعدم التقدير، المقارنة). وعلى مستوى أبعاد تقدير الذات، تبيّن أنّ \emph{البعدين الشخصي
والأسري} كانا الأكثر تأثّرًا (سجّل البعد الأسري تسع حالات بمستوى مرتفع)، يليهما \emph{البعد
الاجتماعي} بتأثّر متوسّط يتفاوت بحسب شدّة الرقابة، ثمّ \emph{البعد الجسمي} بتأثّر ضعيف إلى متوسّط
مع فجوة بين الجنسين. أمّا \emph{البعد الأكاديمي} فكان الأقلّ تأثّرًا، بل تحوّل الضغط الوالدي فيه
لدى بعض المشاركين إلى دافع للإنجاز، لأنّهم فسّروا المتابعة الدراسية بوصفها حرصًا على نجاحهم لا
انتقاصًا من قيمتهم؛ وهي نتيجة تُعدّ من أبرز إسهامات الدراسة، إذ لم تُظهرها البحوث الكمّية التي
تتعامل مع تقدير الذات في صورته العامّة. كما كشفت المقابلات عن مؤشّرات نوعية لم تكن ضمن أبعاد
الدراسة، من أهمّها أثر الصراعات الأسرية، والتناقض بين أقوال الوالدين وممارساتهما، واختلاف إدراك
الإخوة للممارسات الوالدية داخل الأسرة نفسها.

\medskip
\noindent\textbf{مناقشة النتائج في ضوء النظريات والدراسات السابقة (موجزًا).}\ نوقشت هذه النتائج
بعدًا بعدًا؛ فقد فُسِّر تأثّر \emph{البعد الشخصي} في ضوء بومريند (ارتفاع الضبط وضعف الحوار)
وروزنبرغ وكوبر سميث (تشكّل تقدير الذات عبر التقييم المنعكس)، بما يتّفق مع لامبورن وزملائه (1991)
وأيت مولود وإكردوشن بعلي (2018) وشكراوي (2025). وفُسِّر \emph{البعد الأسري} -- الأعلى تأثّرًا --
عبر ماكوبي ومارتن (ضبط مرتفع واستجابة انفعالية منخفضة) وكوبر سميث وماسلو (الحاجة إلى القبول
والتقدير)، متقاطعًا مع الدراسات نفسها. وأُضيء \emph{البعد الاجتماعي} بتصوّر باربر (التحكّم النفسي)
وهارتر (القبول داخل الجماعة)، فظهر تأثّرٌ متوسّط متفاوت. أمّا \emph{البعد الجسمي} فنوقش عبر هارتر
وماسلو مع فجوةٍ بين الجنسين، في حين تميّز \emph{البعد الأكاديمي} بكونه الأقلّ تأثّرًا؛ إذ فسّرت
نتائج الدراسة الحالية -- بخلاف الدراسات الكمّية السابقة التي تعاملت مع تقدير الذات متغيّرًا عامًّا
-- بقاءه محصّنًا لأنّ المراهقين أوّلوا الصرامة الدراسية حرصًا على نجاحهم لا انتقاصًا من قيمتهم،
وهو ما يُعدّ الإضافة النوعية للدراسة. وتباينت جزئيًا الدراسة النيبالية (2023) بحكم خصوصية سياقها
الثقافي.

\medskip
\noindent\textbf{الإجابة عن السؤال المركزي.}\ خلصت الدراسة إلى أنّ علاقة المراهق بتقدير ذاته لا
ترتبط بالممارسات الوالدية في ذاتها وإنّما بالمعنى الذي يمنحه لها؛ ولذلك اختلفت هذه العلاقة من
بعد إلى آخر، ومن مراهق إلى آخر، بل بين الإخوة في الأسرة الواحدة، تبعًا لاختلاف الخبرات وطريقة
تفسيرها. وبذلك تكمن قيمة الدراسة في تقديم قراءة كيفية للعلاقة من منظور المراهق نفسه، وفي إثراء
البحث الكيفي في البيئة المغربية، مع إقرارها بحدودها المنهجية المرتبطة بحجم العيّنة القصدية ووحدة
السياق المؤسّسي.

\medskip
\noindent\textbf{بنية البحث.}\ انتظمت الدراسة في أربعة فصول: خُصّص \textbf{الفصل الأول} للإطار
النظري (المراهقة، ومفهوم الإدراك، وأسلوب المعاملة الوالدية المتسلط وأبعاده، وتقدير الذات وأبعاده،
والتعريفات الإجرائية)؛ وخُصّص \textbf{الفصل الثاني} للإشكالية وأهمّية الدراسة وأهدافها وأسئلتها
وإجراءاتها المنهجية (المنهج، والعيّنة، وأداة جمع المعطيات، ومنهجية تحليل البيانات)؛ وعُرضت في
\textbf{الفصل الثالث} نتائج تحليل المقابلات وفق أبعاد تقدير الذات الخمسة؛ وتناول \textbf{الفصل
الرابع} مناقشة هذه النتائج وتفسيرها في ضوء الإطار النظري والدراسات السابقة، والإجابة عن السؤال
المركزي، واستخلاص أهمّ الاستنتاجات وإبراز الإسهامات والحدود وآفاق البحث المستقبلية.
